\paragraph{A conditional history gives a conditional opening.}
Importing a known book and selecting a possible history answer different
questions. Suppose the disclosed constraints permit closing cash of USD~15.
One compatible synthetic history starts with USD~10 and adds sales of USD~2
and USD~3; another starts with USD~8 and adds sales of USD~3 and USD~4.
Both reach the same cash balance, yet their cumulative-sales memory is USD~5
and USD~7 respectively. A continuation that reads that memory must retain
which history supplied it. Equal closing statements do not erase the difference.

\texttt{ConditionalCompanyOpening/1} accepts an explicit witness for the
typed history and measurement attachment. It replays the complete supplied
constraints before constructing an opening. Every account maps once to a
uniquely identified account with an explicitly declared classification; every
numeric memory field maps once to a nonoverlapping path. Signed amounts,
dimensions and scale are preserved. Missing or duplicate mappings refuse
rather than creating an unallocated balancing amount. In the example, both
openings contain cash of USD~15 and equity credit of USD~15, while their
separate sales memories remain USD~5 and USD~7. A checked inventory example
similarly retains one physical unit separately from inventory of USD~2 and
cash of USD~21.

The boundary between the representations needs explicit premises. This first
adapter selects the final declared end-of-day snapshot and starts the book
on the next day, assuming no unrecorded event across that boundary. Entity
and book must agree, and monetary amounts must already use the same currency
at unit scale. The target context is explicitly synthetic. The source
accounting-basis context and the stated mapping premise remain attached to
it; changing the representation does not establish a basis conversion.
Income and expense accounts must already have zero opening balances.
Otherwise the adapter refuses: it does not invent fiscal closing entries.

The returned provenance retains the original history-and-measurement
definition, selected witness, source snapshot, mapping premises, exclusions
and unresolved obligations. Missing premises yield an unresolved result.
Remaining target obligations or unmapped observations stay visible in a
candidate opening. A valid opening balance report checks the mapped book;
it does not admit the history or authorize a financial continuation. No
investor information cutoff or forecast admissibility has been checked.

This distinction matters for receivables. An aggregate balance can be mapped
exactly without supplying invoice identities, counterparties or due dates.
Those facts must come from an explicit declared or source-supported contract
population before an identified collection method can use them. The adapter
therefore provides a checked conditional starting state while preserving the
remaining method and contract obligations. Other compatible histories remain
represented by the unchanged parent definition; this selected witness neither
enumerates them nor assigns them probabilities.

